\documentclass[10pt,a4paper]{amsart}
\usepackage[margin=19mm]{geometry}
\usepackage{amsmath,amssymb,mathtools}
\usepackage[scr=rsfs,cal=euler]{mathalfa}
\usepackage{xfrac}
\usepackage[unicode,hidelinks,bookmarks=false]{hyperref}
\newcommand{\ie}{i.e.\ }
\newcommand{\cf}{cf.\ }
\newcommand{\resp}{respectively\ }
\newcommand{\Subsection}{\subsection}
\providecommand{\nobreakdash}{\nobreak\hbox{-}}
\setlength{\emergencystretch}{2em}
\multlinegap0pt
\usepackage{bm}
\usepackage{bbm}
\usepackage{dsfont}
\usepackage{xspace}
\usepackage{cancel}
\usepackage{mathtools}
\usepackage{amsthm}
\makeatletter
\@ifundefined{theorem}{\newtheorem{theorem}{Theorem}}{}
\makeatother
\title[Exactness, Cohomology, and Uniqueness in First-Order Differe]{Exactness, Cohomology, and Uniqueness in First-Order Differential Equations}
\author{Hemanta Mandal}
\date{Source: arXiv:2507.16457v2 (CC BY 4.0)}
\begin{document}
\maketitle

\noindent\emph{Source and licence.} Exactness, Cohomology, and Uniqueness in First-Order Differential Equations. Version arXiv:2507.16457v2 (CC BY 4.0), distributed under the Creative Commons Attribution 4.0 International licence, as recorded in the arXiv licence metadata for this exact version and shown on the abstract page https://arxiv.org/abs/2507.16457v2. Full rights notice: licenses/arxiv-2507.16457-CC-BY-4.0.txt.\par\smallskip

\noindent\textit{Editorial scope.} Counted unit: the Theorem whose source label is \texttt{None} in the article (source0.tex lines 182--206, source label \texttt{None}), delivered together with the article's own complete original proof of it (source0.tex lines 208--224, one \texttt{proof} environment of 17 lines). No mathematics is written, repaired or re-derived: statement and proof are the article's own text line for line. Required context retained uncounted: none. Every definition, display and label of those lines is retained unchanged. Omitted from this package and not redistributed: 1--181, 207--207, 225--422. Nothing inside a retained excerpt is omitted or altered: every retained body is delivered exactly as the article's lines are written, so no omission receipt of any kind accompanies this package. Citations: the retained excerpts carry no citation command at all, so no bibliography entry of the article is transcribed and no reference is redistributed. Reference adaptation: none was needed and none was made; every reference inside the delivered unit resolves inside it, so no unresolved reference or citation is produced. Printed slips kept literal: no printed word, symbol, hypothesis or number of the counted statement or of its proof was altered. Layout and class: the article's own class is \texttt{amsart} with packages including graphicx, amssymb, booktabs, pifont; the wrapper is the lane's pinned \texttt{amsart} editorial wrapper at 10pt on A4, which restates the theorem-like environments the retained text needs and adds the article's own macro definitions that the retained text needs (none), with extra wrapper packages (bm, bbm, dsfont, xspace, cancel, mathtools). The delivered numbering is the unit's own. Assets: no retained span contains a figure environment, a graphics-inclusion command, a PDF-page inclusion, a table or a TikZ picture; no image file is copied into this package, the package declares no asset dependency and no image file is redistributed with it. Licence: version arXiv:2507.16457v2 of this article is distributed under the Creative Commons Attribution 4.0 International licence, as recorded in the arXiv licence metadata for this exact version and shown on the abstract page https://arxiv.org/abs/2507.16457v2. A full rights notice is delivered at licenses/arxiv-2507.16457-CC-BY-4.0.txt. Characters: every retained excerpt is pure ASCII, so the delivered engine, XeLaTeX with its default Latin Modern text font, reports no missing character.
\begin{theorem}[Global Solvability via Trivial de\,Rham Class]
Let \( U \subset \mathbb{R}^2 \) be an open, connected (not necessarily simply connected) domain, and let
\[
\omega = M(x, y)\, dx + N(x, y)\, dy
\]
be a smooth, closed 1-form on \( U \), i.e.,
\[
\frac{\partial N}{\partial x} = \frac{\partial M}{\partial y}.
\]
Suppose that the de\,Rham cohomology class \( [\omega] \in H^1_{\mathrm{dR}}(U) \) is trivial (i.e., \( \omega \) is exact).

Then there exists a smooth function \( F : U \to \mathbb{R} \) such that
\[
dF = \omega,
\]
and the general solution to the differential equation
\[
M(x, y)\, dx + N(x, y)\, dy = 0
\]
is given by
\[
F(x, y) = C,
\]
for arbitrary constants \( C \in \mathbb{R} \). Moreover, such a function \( F \) is unique up to an additive constant.
\end{theorem}

\begin{proof}
Since \( \omega \) is a closed 1-form on \( U \), it defines a cohomology class in \( H^1_{\mathrm{dR}}(U) \). The hypothesis \( [\omega] = 0 \) means that \( \omega \) is exact; that is, there exists a smooth function \( F : U \to \mathbb{R} \) such that
\[
\omega = dF.
\]
Substituting this into the differential equation, we have
\[
dF = 0 \quad \Rightarrow \quad F(x, y) = C,
\]
where \( C \) is a constant. Thus, the solution curves lie along level sets of \( F \), which globally solves the equation.

To show uniqueness: if \( \tilde{F} \) is another function satisfying \( d\tilde{F} = \omega \), then
\[
d(F - \tilde{F}) = dF - d\tilde{F} = \omega - \omega = 0,
\]
so \( F - \tilde{F} \) is constant on each connected component of \( U \). Since \( U \) is connected, \( F - \tilde{F} \) must be a global constant. Therefore, \( F \) is unique up to an additive constant.
\end{proof}

\end{document}
